\section{Movement in concrete affine formulations}
\label{sec:concrete-movement}

We compute distances in the restricted barrier metric for grouped
balls, norm trees, balance slices, and shared spectral cones; the number
of cones in the ambient product does not determine this distance.  The
barrier-height argument and the grouped, packed, and unweighted tree
central paths also appear in the companion
manuscript~\cite{CentralPathCompanion}; we include the proofs to make
the comparisons self-contained.  Section~\ref{sec:tree-distance} gives
the exact distance coefficient for norm trees with a fixed objective.
Throughout, distance means the infimum of lengths of piecewise $C^1$
curves in the indicated open affine slice, with its displayed Hessian
metric; the distance to a target set is the infimum over its points.
We write $[u]_+=\max\{u,0\}$.

\subsection{Barrier height and simultaneous slack contraction}

\begin{lemma}[Two endpoint estimates]\label{lem:cm-height}
Let $F$ be a barrier whose restricted gradient satisfies
$\norm{DF}_{F,*}\leq\sqrt\nu$.  For every reference $x^0$ and target
set $E$,
\begin{equation}\label{eq:cm-height}
 d_F(x^0,E)\geq
 \frac{[\inf_{y\in E}F(y)-F(x^0)]_+}{\sqrt\nu}.
\end{equation}
If, in addition, $F=-\sum_{i=1}^H\log q_i$ for positive $C^2$
concave functions on a convex domain and the Hessian is positive
definite, then
\begin{equation}\label{eq:cm-log-slack}
 d_F(x,y)\geq
 \left(\sum_i\left|\log\frac{q_i(y)}{q_i(x)}\right|^2\right)^{1/2}.
\end{equation}
In particular, if $\prod_iq_i(y)\leq(A\epsilon/H)^H$ throughout
$E_\epsilon$ and $\bar q_0=(\prod_iq_i(x^0))^{1/H}$, then
\[
 d_F(x^0,E_\epsilon)\geq
 \sqrt H\left[\log\frac{H\bar q_0}{A\epsilon}\right]_+.
\]
\end{lemma}
\begin{proof}
Integrate $|DF[\dot x]|\leq\sqrt\nu\norm{\dot x}_F$ along a
curve for the first assertion.  Concavity gives
\[
 D^2(-\log q_i)[u,u]
 =\frac{Dq_i[u]^2}{q_i^2}-\frac{D^2q_i[u,u]}{q_i}
 \geq (D\log q_i[u])^2.
\]
Thus the map $x\mapsto(\log q_i(x))_i$ into Euclidean space does not
increase lengths.
Integration and Cauchy--Schwarz prove the remaining assertions.
\end{proof}

Both estimates concern the whole accurate set, not the length of a
particular central path.  We record the resulting bound on bounded moves.  A feasible chord of starting
local norm $r<1$ has length at most $-\log(1-r)$, by integrating the
self-concordant Hessian comparison.  Hence if a trajectory from $x^0$,
with $\norm{x^0-\bar x}_{\bar x}\leq\rho<1$, reaches $E_\epsilon$ by
chords of starting local norms $r_j<1$, and $D_\epsilon$ is a lower
bound on $d_F(\bar x,E_\epsilon)$, then
\begin{equation}\label{eq:cm-movement-budget}
 \sum_j-\log(1-r_j)\geq
 [D_\epsilon+\log(1-\rho)]_+.
\end{equation}
In particular, for an integer $m\geq1$ and $0<R<1$, rounds of at most
$m$ chords of starting local norm at most $R$ require at least $[D_\epsilon+\log(1-\rho)]_+/[-m\log(1-R)]$ rounds.
With $M$ chords in total, where $M\geq1$ is an
integer, some $r_j$ is at least
$1-\exp(-[D_\epsilon+\log(1-\rho)]_+/M)$.
These statements restrict only the geometry of the primal trajectory
and allow arbitrary computation between the counted moves.

\subsection{Grouped balls and unrestricted PSD completion}

For each of $k$ balls, partition the coordinates into $h$ nonempty groups,
and use the grouped domain from~\eqref{eq:grouped-domain-concrete}:
\[
 q_{aG}=s_{aG}-\norm{w_{aG}}^2>0,\qquad
 \sum_Gs_{aG}=1,\qquad F=-\sum_{a,G}\log q_{aG}.
\]
Choose unit directions $c_a$, maximize
$\ell=\sum_a\ip{c_a}{w_a}$, and put $H=kh$.
The optimum is $k$.  A PSD packing puts these $H$ group columns into
arbitrary real matrix blocks $W_j$, pads shorter columns by zeros, and
replaces the scalar residuals by
\[
 D_j=S_j-W_j^TW_j\succ0,\qquad F=-\sum_j\log\det D_j,
 \qquad \sum_{G\text{ of }a}(S_{j(a,G)})_{GG}=1.
\]
All off-diagonal entries of $S_j$ are free completion variables.  Each column occurs once
and contains precisely its assigned group coordinates.  This includes
the one-factor packing of Section~\ref{sec:concrete-barriers}.

\begin{proposition}[Completion variables do not shorten the distance]
\label{prop:cm-grouped}
For both formulations the analytic center has $w=0$ and all residual
eigenvalues equal to $1/h$.  The central path minimizing $F-\tau\ell$
is given by
\begin{equation}\label{eq:cm-radial-profile}
 D(\tau)=\sqrt{h^2+\tau^2},\qquad
 q(\tau)=\frac{2}{D(\tau)+h},\qquad
 r(\tau)=\frac{\tau}{D(\tau)+h},
\end{equation}
with $w_a=r(\tau)c_a$ and $q_{aG}=q(\tau)$ or
$D_j=q(\tau)I$.  From every central reference $z(\tau_0)$,
\begin{equation}\label{eq:cm-grouped-distance}
 d_F\bigl(z(\tau_0),\{k-\ell\leq\epsilon\}\bigr)
 \geq\sqrt H\left[\log\frac{Hq(\tau_0)}{2\epsilon}\right]_+.
\end{equation}
The bound allows arbitrary motion in the completion variables.
\end{proposition}
\begin{proof}
For the grouped barrier, allocation stationarity makes all slacks of
one ball equal, and leaf stationarity gives $2w_{aG}/q=\tau c_{aG}$.
The root allocation identity is $hq+\norm{w_a}^2=1$.
Solving it gives~\eqref{eq:cm-radial-profile}.
For packing, stationarity in off-diagonal entries of $S_j$ makes
$D_j^{-1}$ diagonal; diagonal allocation stationarity equates its
entries within each ball.  The remaining equations are the same.
Strict convexity on the stated affine coordinates gives uniqueness.

Let $\delta_a=1-\ip{c_a}{w_a}>0$.  In either representation the sum
of residual diagonal entries belonging to ball $a$ is
$1-\norm{w_a}^2\leq2\delta_a$.  Hadamard's determinant inequality
for packing, followed by AM--GM within and across balls, gives
\[
 \prod_{a,G}q_{aG}\leq(2\epsilon/H)^H,
 \quad\hbox{or}\quad
 \prod_j\det D_j\leq(2\epsilon/H)^H
 \qquad\text{if }\sum_a\delta_a\leq\epsilon.
\]
The grouped barrier has parameter $H$.  So does the unrestricted
matrix epigraph barrier before allocation restriction: the shear
calculation in the proof of Theorem~\ref{thm:packing-intrinsic} gives
squared gradient norm equal to the number of residual columns, and
restriction can only decrease it.  At $z(\tau_0)$ the residual
determinant product is $q(\tau_0)^H$.
Lemma~\ref{lem:cm-height} proves the result, with no constraint on
the completion variables along the path.
\end{proof}

A blockwise version of this estimate retains more information than the
total barrier height.  If $c_j$ is the number of columns in block
$j$, differentiation gives
\[
 D^2F_j[(\dot S_j,\dot W_j),(\dot S_j,\dot W_j)]
 =\tr\bigl((D_j^{-1/2}\dot D_jD_j^{-1/2})^2\bigr)
   +2\tr(D_j^{-1}\dot W_j^T\dot W_j),
 \quad
 \dot D_j=\dot S_j-\dot W_j^TW_j-W_j^T\dot W_j.
\]
Consequently every path from $x$ to $y$ has length at least
\begin{equation}\label{eq:cm-packed-log-vector}
 \left(\sum_j\frac1{c_j}
 \left|\log\frac{\det D_j(y)}{\det D_j(x)}\right|^2\right)^{1/2}.
\end{equation}
Indeed, trace Cauchy--Schwarz bounds each squared logarithmic
derivative by $c_j$ times its block metric, and integration gives
the assertion.  No off-diagonal entry is fixed in this argument.

For different group counts $h_a\geq1$ and positive objective weights
$\lambda_a$, set $H=\sum_ah_a$, $p_a=h_a/H$, and
$\Delta_{\mathrm{eff}}=\prod_a(\lambda_a/p_a)^{p_a}$.
In both grouped and packed formulations, the residual diagonal entries
of source $a$ at the analytic center $\bar z$ equal $1/h_a$, the
off-diagonal residuals are zero, and
\begin{equation}\label{eq:cm-packed-entropy}
 d_F\left(\bar z,
 \left\{\sum_a\lambda_a(1-\ip{c_a}{w_a})\leq\epsilon\right\}\right)
 \geq\sqrt H\left[\log\frac{\Delta_{\mathrm{eff}}}{2\epsilon}\right]_+.
\end{equation}
To verify it, put $e_a=\lambda_a(1-\ip{c_a}{w_a})$.
Hadamard's inequality and AM--GM within each source bound the residual determinant product
by $\prod_a(2e_a/(\lambda_ah_a))^{h_a}$; its maximum for
$\sum_ae_a\leq\epsilon$ occurs at $e_a=\epsilon p_a$.
Compare with the center product $\prod_ah_a^{-h_a}$ and use
Lemma~\ref{lem:cm-height}.  In particular, the one-block lift of
order $s+k$ with one length-$s$ column per ball has $H=k$, despite the
unrestricted completion variables in its single PSD factor.
For the one-column-per-ball Hermitian packing of
\eqref{eq:packing-domain-concrete}, these endpoint and central-path
statements also hold over $\C$ and $\HH$, with arbitrary nonzero real
column subspaces and unit support directions $c_a$ in them.  Use the real trace
and, over $\HH$, the Jordan determinant.  Hermitian Hadamard gives the
same diagonal product bound; Theorem~\ref{thm:packing-intrinsic}
gives parameter $k$; and stationarity in every real component of the
free off-diagonal entries forces $D^{-1}$ to be diagonal.  Column
stationarity in its real subspace then gives $w_a=rc_a$, as above.
Thus the arc and tube formulas below apply to this Hermitian
specialization with $h=1$ as well.

\subsection{Norm trees: uniform bounds and weighted exact central paths}

Let $T$ be a finite rooted tree whose internal nodes have at least two
children.  Its leaves index the coordinates of a Euclidean unit ball.
Write $I(v)$ and $J(v)$ for the internal and leaf children of $v$, respectively, and set
\begin{equation}\label{eq:cm-tree-domain}
 q_v=t_v^2-\sum_{u\in I(v)}t_u^2-\sum_{j\in J(v)}w_j^2>0,
 \qquad t_v>0,\qquad t_{\mathrm{root}}=1.
\end{equation}
All blocks must be on the positive Lorentz sheet.  Assign fixed weights
$\omega_v\geq1$ and use $F_T=-\sum_v\omega_v\log q_v$.
Put $W=\sum_v\omega_v$, and let $\Omega_v$ be the sum of the
weights of internal nodes in the subtree rooted at $v$, including $v$.
For a unit leaf direction $c$, put
$\alpha_v=\norm{c|_{\text{leaves below }v}}$, so
$\alpha_{\mathrm{root}}=1$.

\begin{proposition}[Weighted central profile]\label{prop:cm-tree-central}
The unique minimizer of $F_T-\tau\ip{c}{w}$ has
\begin{equation}\label{eq:cm-weighted-tree-central}
 \begin{split}
 q_v&=\omega_v q,\qquad
 q=\frac{2}{\sqrt{W^2+\tau^2}+W},\qquad
 r=\frac{\tau}{\sqrt{W^2+\tau^2}+W},\\
 w&=rc,\qquad t_v^2=\Omega_v q+r^2\alpha_v^2.
 \end{split}
\end{equation}
For $\tau>0$ its squared speed per unit $\log\tau$ is
\begin{equation}\label{eq:cm-weighted-tree-speed}
 \left\|\frac{dz}{d\log\tau}\right\|_{F_T}^2
 =W\left(1-\frac{W}{\sqrt{W^2+\tau^2}}\right).
\end{equation}
The formulas allow zero leaf coefficients and zero subtree norms.
\end{proposition}
\begin{proof}
Internal-coordinate stationarity is
$-2\omega_vt_v/q_v+2\omega_{p(v)}t_v/q_{p(v)}=0$.
Since $t_v>0$, all $q_v/\omega_v$ are equal.  Leaf stationarity is
$2w_j/q=\tau c_j$.  The determinant identities telescope to
$\sum_vq_v=1-\norm w^2$, giving $Wq+r^2=1$ and the stated solution.
The subtree identities determine the positive $t_v$ uniquely.
The weighted barrier is strictly convex on the slice $t_{\mathrm{root}}=1$, so the stationary point is the unique minimizer.  Differentiating
restricted central stationarity gives $H_\tau z'(\tau)=\nabla\ell$;
therefore squared logarithmic speed is $\tau^2r'(\tau)$, which is
\eqref{eq:cm-weighted-tree-speed}.
\end{proof}

For $k$ independent trees with the same total weight $W$, unit
objective weights, and $L_\omega=kW$, all profiles agree and their
squared speeds add.
The same expression holds for grouped and packed balls with $W=h$
and $L_\omega=H$.  Since restricted central stationarity gives
$\nabla F=\tau\nabla\ell$, their squared restricted gradient norm
along the central path tends to $H$.  Thus the standard grouped and
packed barriers have exact parameter $H$, not just the upper bound
used in Proposition~\ref{prop:cm-grouped}.  Define
\begin{equation}\label{eq:cm-arc-function}
 \mathcal A(r)=\sqrt2\left[
 \sqrt2\operatorname{artanh}\frac{\sqrt2r}{\sqrt{1+r^2}}
 -\operatorname{artanh}\frac{r}{\sqrt{1+r^2}}\right],
 \qquad 0\leq r<1.
\end{equation}
The exact central length from the analytic center to profile $r$ is
\begin{equation}\label{eq:cm-exact-arc}
 \sqrt{L_\omega}\,\mathcal A(r),\qquad
 \mathcal A(r)=\log\frac1{1-r}+O(1)\quad(r\uparrow1).
\end{equation}
Indeed, $\tau=2Wr/(1-r^2)$ transforms the length element into
$\sqrt{2L_\omega}\sqrt{1+r^2}\,dr/(1-r^2)$; the substitution
$y=r/\sqrt{1+r^2}$ gives~\eqref{eq:cm-arc-function}.
For $0<\epsilon\leq k$, the central path to gap $\epsilon$ uses
$r=1-\epsilon/k$. At $\epsilon=k$ its length is zero; larger tolerances
are already met at the analytic center.
An arc of length $a$ ends within local norm $e^a-1$ of its start,
measured at the start: after arclength $s$, self-concordance bounds the
ratio of local norms by $e^s$, and integration gives the estimate.
Partitioning the central path into arcs of length
$\log(1+R)$, where $0<R<1$, therefore gives at most
\begin{equation}\label{eq:cm-arc-chords}
 \left\lceil
 \frac{\sqrt{L_\omega}\,\mathcal A(1-\epsilon/k)}{\log(1+R)}
 \right\rceil
\end{equation}
feasible chords of starting local norm at most $R$.

The next estimate holds uniformly at finite accuracy and allows
different tree weights and objective weights.

\begin{proposition}[Entropy scale for different trees]
\label{prop:cm-tree-entropy}
For source $a$, let $W_a=\sum_v\omega_{av}$ and let
$\lambda_a>0$ multiply its unit support objective.  Put
\[
 L_\omega=\sum_aW_a,\quad p_a=W_a/L_\omega,\quad
 \Delta_{\mathrm{eff}}=\prod_a(\lambda_a/p_a)^{p_a}.
\]
If the displayed restricted barrier
$F=-\sum_{a,v}\omega_{av}\log q_{av}$ satisfies
$\norm{DF}_{F,*}\leq\sqrt\nu$, then from its analytic center $\bar z$,
\begin{equation}\label{eq:cm-tree-entropy}
 d_F\left(\bar z,
 \left\{\sum_a\lambda_a(1-\ip{c_a}{w_a})\leq\epsilon\right\}\right)
 \geq\frac{L_\omega}{\sqrt\nu}
 \left[\log\frac{\Delta_{\mathrm{eff}}}{2\epsilon}\right]_+.
\end{equation}
The product-cone value $\nu=2L_\omega$ is always valid.
For unit barrier weights, write $b_a$ for the number of internal
nodes.  One may use the exact parameter
$\nu=\sum_a\nu_{T_a}$ from Theorem~\ref{thm:tree-parameter}, where
$\nu_{T_a}=2b_a-2$ if the root has only internal children and
$\nu_{T_a}=2b_a-1$ otherwise.
\end{proposition}
\begin{proof}
At the center $q_{av}^0=\omega_{av}/W_a$.  Write
$e_a=\lambda_a(1-\ip{c_a}{w_a})$.  Telescoping gives
$\sum_vq_{av}=1-\norm{w_a}^2\leq2e_a/\lambda_a$.
Weighted AM--GM within a tree gives
\[
 \prod_v\left(\frac{q_{av}}{q_{av}^0}\right)^{\omega_{av}}
 \leq(2e_a/\lambda_a)^{W_a}.
\]
Under $\sum_ae_a\leq\epsilon$, the product on the right is
maximized at $e_a=\epsilon p_a$.  Thus
$F(y)-F(\bar z)\geq L_\omega\log(\Delta_{\mathrm{eff}}/(2\epsilon))$.
Apply Lemma~\ref{lem:cm-height}.
\end{proof}

For unit objective weights, $\Delta_{\mathrm{eff}}=\exp(H(p))$,
where $H(p)=-\sum_ap_a\log p_a$ is Shannon entropy in natural units.
In general $\Delta_{\mathrm{eff}}\leq\sum_a\lambda_a$, with equality
exactly when the objective weights are proportional to $W_a$.
For $k$ unweighted identical trees with $b$ internal nodes,
$L=L_\omega=kb$ and the bound is
$(L/\sqrt\nu)[\log(k/(2\epsilon))]_+$.
Since $L\leq\nu<2L$, this bound and~\eqref{eq:cm-arc-chords} show
that reaching the target takes
$\Theta_R(\sqrt L\log(k/\epsilon))$ chords of starting local norm at
most $R$ for $0<\epsilon\leq k/4$, with
constants independent of the tree.
The lower estimate applies to arbitrary feasible paths.
Section~\ref{sec:tree-distance} determines the exact leading coefficient
for a fixed objective, which can be smaller than the central-path
coefficient.

\subsection{A separate control on central-path parameter changes}

The endpoint bounds above hold whether or not accepted iterates stay
near the central path.  If they do, each increase of the path parameter
is also bounded.  This refines the short-step analysis for grouped
balls, and the proof applies to the tree profile as well.

\begin{proposition}[Tube parameter increments]\label{prop:cm-tube}
Consider a profile above with $k$ sources of unit objective weight and
common scalar parameter $W\geq1$ (grouped or packed balls with $W=h$,
or trees of total weight $W$), $L_\omega=kW$, and objective optimum $k$.
Fix $a>0$, $0<R<1$, an integer $m\geq1$, and
\[
 \gamma_a=1+a-\sqrt{1+a^2},\qquad
 b_a=\sqrt{1-(1+a^2)^{-1/2}},\qquad 0\leq\rho<\gamma_a.
\]
Suppose $aW\leq\tau_0\leq\cdots\leq\tau_T$, accepted points satisfy
$\norm{x_j-z(\tau_j)}_{z(\tau_j)}\leq\rho$, and each round, from
$x_j$ to $x_{j+1}$, uses at most $m$ feasible chords of starting norm at
most $R$.
Put
\[
 B_\rho=\frac{\rho}{(1-\rho)^2},\qquad
 A_{\rho,R,m}=(1-R)^{-m}-1+
 B_\rho\bigl(1+(1-R)^{-m}\bigr).
\]
Then every round satisfies
\begin{equation}\label{eq:cm-tube-increment}
 \log\frac{\tau_{j+1}}{\tau_j}
 \leq\frac{A_{\rho,R,m}}{(1-\rho)b_a\sqrt{L_\omega}}.
\end{equation}
If the final objective gap is at most $\epsilon$, and
$\Delta=L_\omega/\tau_0$, then
\begin{equation}\label{eq:cm-tube-count}
 T\geq\frac{(1-\rho)b_a\sqrt{L_\omega}}{A_{\rho,R,m}}
 \left[\log\frac{(\gamma_a-\rho)\Delta}{\epsilon}\right]_+.
\end{equation}
\end{proposition}
\begin{proof}
All covectors are restricted to the affine hull.  Hessian comparison
and integration from $z(\tau_j)$ to $x_j$ give
\[
 \norm{\nabla F(x_j)-\tau_j\nabla\ell}_{x_j,*}\leq B_\rho.
\]
For one chord of starting norm at most $R$, the gradient change in its starting dual norm is
at most $R/(1-R)$; transporting a dual norm backwards across that
chord costs at most $(1-R)^{-1}$.  A geometric sum over at most $m$
chords, including both endpoint residuals, consequently gives
\[
 (\tau_{j+1}-\tau_j)\norm{\nabla\ell}_{x_j,*}
 \leq A_{\rho,R,m}.
\]
The exact speed formula and the tube comparison give
$\tau_j\norm{\nabla\ell}_{x_j,*}
\geq(1-\rho)b_a\sqrt{L_\omega}$.
Use $\log u\leq u-1$ to prove~\eqref{eq:cm-tube-increment}.
For $u=\tau/W$, the central gap $g$ satisfies
$\tau g/L_\omega=1+u-\sqrt{1+u^2}\geq\gamma_a$.
The objective change from a center to a point within local norm $\rho$
of it, measured at the center, is
at most $\rho\sqrt{L_\omega}/\tau\leq\rho L_\omega/\tau$.
Thus accuracy forces
$\tau_T\geq(\gamma_a-\rho)L_\omega/\epsilon$.
Sum~\eqref{eq:cm-tube-increment} to finish.
\end{proof}

\subsection{The balance slice attains the scale set by its certificate rank}

Return to $Z$ in~\eqref{eq:balance-slice}, with total Jordan rank
$\rho=\sum_i r_i$, and use the restricted standard barrier
$\Phi=-\sum_i\log\det X_i$.
Choose the frame primitive $c=c_{11}$ and put
$S=(e_1-c,e_2,\ldots,e_L)$ and $Q=\rho-1$.
The gap $g(X)=\ip{S}{X}=1-\ip{c}{X_1}$ vanishes only at
$(c,0,\ldots,0)$: positivity and vanishing pairing force every block
into the corresponding complementary face, and trace one fixes its
coefficient.  This point is feasible because all balance moments
vanish on the chosen frame.

\begin{proposition}[Exact leading distance on the balance slice]
\label{prop:cm-balance}
Let $X_i^0=e_i/\rho$ and $g_0=Q/\rho$.  For $0<\epsilon<g_0$,
\begin{equation}\label{eq:cm-balance-distance}
 \sqrt Q\log\frac{g_0}{\epsilon}
 \leq d_\Phi(X^0,\{g\leq\epsilon\})
 \leq\sqrt Q\log\frac{g_0}{\epsilon}
       +\log\bigl(\rho(1-\epsilon)\bigr).
\end{equation}
Thus, for this fixed instance, the leading coefficient of the distance
in $\log(1/\epsilon)$ is exactly $\sqrt{\rho-1}$.
\end{proposition}
\begin{proof}
The certificate $S$ has total Jordan rank $Q$, determinant one on its
support, and exposed determinant $\rho^{-Q}$ at $X^0$.
Theorem~\ref{thm:movement-minor} therefore gives scale
$\Delta=Q/\rho=g_0$ and the lower bound.  More explicitly, the
negative logarithm of the exposed determinant has dual gradient norm
at most $\sqrt Q$ on the slice, and spectral AM--GM bounds that
determinant at gap $\epsilon$ by $(\epsilon/Q)^Q$.

The central path minimizing $\Phi+\tau g$ has
\[
 X_1=ac+b(e_1-c),\qquad X_i=be_i\ (i>1),\qquad
 a+Qb=1,\quad \frac1b=\tau+\frac1a.
\]
This point is stationary on the whole slice: its balance moments
vanish, and $\nabla\Phi+\tau S$ is the total-trace normal $-e/a$.  Strict convexity gives uniqueness and also proves the
center assertion at $\tau=0$.  Writing $g=Qb$ gives
\[
 \tau=\frac{Q-(Q+1)g}{g(1-g)},\qquad
 g(0)=g_0,\qquad g(\tau)\sim Q/\tau.
\]
The frame-diagonal metric has central length
\[
 \int_\epsilon^{g_0}
 \sqrt{\frac Q{g^2}+\frac1{(1-g)^2}}\,dg.
\]
Bound its integrand above by $\sqrt Q/g+1/(1-g)$ to obtain
the upper bound in~\eqref{eq:cm-balance-distance}.
\end{proof}

Theorem~\ref{thm:balance-slices} independently proves that every
barrier on this body has parameter at least $\rho-1$.  The distance
statement above concerns only the displayed standard barrier.  Its
squared coefficient equals this parameter bound, but that does not
prove the same distance bound for every optimal coupled barrier.

\subsection{Sharing spectral cones does not change the reduced movement}

Use the real or complex matrix balls and orientations of
Section~\ref{sec:spectral-chordal}, with $r_a\leq c_a$, and put
$R_0=\sum_ar_a$.  Their common reduced barrier is
\[
 \Phi(X)=-\sum_a\log\det(I_{r_a}-X_aX_a^*).
\]
Let $C_a=(I_{r_a}\ 0)$, and maximize
$\ell(X)=\sum_a\operatorname{Re}\tr(C_a^*X_a)$, whose optimum
is $R_0$.

\begin{proposition}[Grouping-independent spectral distance]
\label{prop:cm-spectral}
For every shared-cone grouping and its fixed-identity chordal block-star
realization,
\begin{equation}\label{eq:cm-spectral-distance}
 d_\Phi(0,\{R_0-\ell\leq\epsilon\})
 \geq\sqrt{R_0}\left[\log\frac{R_0}{2\epsilon}\right]_+.
\end{equation}
The exact central path is $X_a(\tau)=r(\tau)C_a$, where
$r(\tau)=\tau/(\sqrt{1+\tau^2}+1)$.  Its length to gap
$0<\epsilon<R_0$ is
$\sqrt{R_0}\mathcal A(1-\epsilon/R_0)$.
Consequently, for this fixed instance, the leading coefficient of the
distance in $\log(1/\epsilon)$ is exactly $\sqrt{R_0}$.
\end{proposition}
\begin{proof}
Let $\sigma_{ai}<1$ be the singular values of $X_a$.
Von Neumann's trace inequality gives
$\sum_{a,i}(1-\sigma_{ai})\leq R_0-\ell(X)$.
Since $1+\sigma_{ai}\leq2$, AM--GM gives
\[
 \prod_a\det(I-X_aX_a^*)
 \leq(2\epsilon/R_0)^{R_0}
 \qquad\text{when }R_0-\ell(X)\leq\epsilon.
\]
Theorem~\ref{thm:shared-spectral-parameter} gives the restricted
parameter $R_0$, and $\Phi(0)=0$; apply
Lemma~\ref{lem:cm-height}.  The barrier gradient at $rC_a$ is
$2rC_a/(1-r^2)=\tau C_a$, so the path is stationary, hence optimal by
convexity, in all variables.
Along the path the barrier is $-R_0\log(1-r^2)$, and its length
is exactly the scalar arc already computed.
The two bounds have the same leading coefficient.
Finally, Proposition~\ref{prop:block-star} identifies the restricted
barrier literally with $\Phi$, so the metric and every displayed
distance and path formula are identical in those formulations.
\end{proof}

For vector rows these bodies are products of Euclidean balls.
Maximal sharing can therefore reduce the number of cones without
changing this reduced path geometry.
Equations~\eqref{eq:cm-movement-budget} and
\eqref{eq:cm-spectral-distance} give the same bounded-chord lower
bound for all these descriptions.  This conclusion depends only on
their common barrier after affine elimination; it gives neither a query
lower bound nor the cost of evaluating the ambient Newton system.
